\documentclass[dvipdfmx,twoside]{jsarticle}
\usepackage{amsmath,amssymb}
\usepackage{CJKutf8}
\usepackage{okumacro}
\usepackage{booktabs}
\usepackage{array}
\usepackage{xcolor}
\usepackage{colortbl}
\usepackage{tipa}
\usepackage{tikz}
\usepackage{fancybox}
\usepackage{wasysym}
\usepackage{pifont}
\usepackage{graphicx}
\usepackage{wrapfig}
\usepackage{geometry}
\geometry{margin=2cm}
\usepackage{fancyhdr}
\usepackage[utf8]{inputenc}
\usepackage{CJK}
% 设置页面样式
\pagestyle{fancy}
\fancyhf{}
\renewcommand{\headrulewidth}{0pt}
\fancyhead[RO]{数学–\thepage}
\fancyhead[LE]{数学–\thepage}

% 重新定义 plain 样式
\fancypagestyle{plain}{%
  \fancyhf{}%
  \renewcommand{\headrulewidth}{0pt}%
  \fancyhead[RO]{数学-\thepage}%
  \fancyhead[LE]{数学-\thepage}%
}
% 自定义A方框命令
\newcommand{\abb}[1]{%
\begin{tikzpicture}[baseline]
\node[draw=black, 
      rectangle, 
      minimum width=0.8cm, 
      minimum height=0.3cm, 
      fill=gray!25, 
      font=\bfseries,
      line width=1pt,
      inner sep=2pt,
      anchor=base] {#1};
\end{tikzpicture}%
}
\newcommand{\ab}[1]{%
\begin{tikzpicture}[baseline]
\node[draw=black, 
      rectangle, 
      minimum width=0.8cm, 
      minimum height=0.3cm, 
      font=\bfseries,
      line width=1pt,
      inner sep=2pt,
      anchor=base] {#1};
\end{tikzpicture}%
}
\newcommand{\sab}[1]{%
\begin{tikzpicture}[baseline]
\node[draw=black, 
      rectangle, 
      minimum width=0.4cm, 
      minimum height=0.05cm, 
      font=\bfseries,
      line width=1pt,
      inner sep=0.5pt,
      anchor=base] {#1};
\end{tikzpicture}%
}

\newcommand{\maru}[1]{\tikz[baseline=-0.7ex]{
    \node[shape=circle,draw,inner sep=1pt,minimum size=5pt,anchor=center] {\footnotesize #1};}}
\definecolor{headercolor}{RGB}{220,220,220}
\definecolor{rowcolor1}{RGB}{245,245,245}
\definecolor{rowcolor2}{RGB}{255,255,255}

% \title{\vspace{-1.5cm} 2025年8月羚課文科数学月考 }
% \author{\textnosfal{Linc\ -\ 伊}}
\date{}
\begin{document}

\begin{CJK}{UTF8}{ipxm}  % 使用ipxm字体
  %月考表纸部分
\begin{center}

\vspace*{5cm}

% 学校logo
\includegraphics[width=5cm]{pics/1.jpg}

\vspace{2cm}

% 主标题
{\fontsize{24}{30}\selectfont\bfseries\sffamily
2025年10月\\
\vspace{1em}
羚課理科数学月考
}

\end{center}
\newpage
%第一部分
\noindent
\begin{tikzpicture}
\node[draw, rectangle, minimum width=1cm, minimum height=0.8cm, font=\Huge] {I};
\end{tikzpicture}
\\
\\
\textbf{問1}\qquad 2つの2次関数

\[
  f(x) = \frac{1}{3}x^2 + ax - b, \qquad g(x) = -x^2 + cx + b
\]

が、次の2つの条件 (A), (B) を満たすような $a$, $b$, $c$ を求めよう。\\[0.5em]

\begin{itemize}
  \item[(A)] グラフ $y=f(x)$ とグラフ $y=g(x)$ は2つの直線 $x=-1,~x=2$ 上で交わっている。\\[0.5em]
  \item[(B)] $g(x)$ の最大値と $f(x)$ の最小値の差は $\dfrac{19}{6}$ である。
\end{itemize}
\ \\[0.5em]
条件(A) より2つの2次関数のグラフは直線 $x=-1$ 上で交わっているから\\[0.5em]
\[
  3a + \ab{\color{red}{\textsf{A}}} b - \ab{\color{red}{\textsf{B}}} c = \ab{\color{red}{\textsf{C}}} \qquad \cdots\ \maru{1}
\]
であり、また、直線 $x = 2$ 上で交わっているから\\[0.5em]
\[
  -\ab{\color{red}{\textsf{D}}} a +3b+\ab{\color{red}{\textsf{E}}} c = 8\qquad \cdots\ \maru{2}
\]
を得る。①と②より、$b = \dfrac{\ab{\color{red}{\textsf{F}}}}{\ab{\color{red}{\textsf{G}}}}$ である。\\[0.5em]

\noindent
次に、条件(B) と $b = \dfrac{\abb{\color{red}{\textsf{F}}}}{\abb{\color{red}{\textsf{G}}}}$ より\\[0.5em]
\[
  \ab{\color{red}{\textsf{H}}} a^2 + c^2 = \ab{\color{red}{\textsf{I}}}
\]
を得る。
\newpage
\begin{center}
- 計算欄 (memo) -
\end{center}
\newpage

\noindent
\textbf{問2}\qquad 2つの袋A、Bがある。Aの袋には白カードが4個，赤カードが1個入っており，Bの袋には白カードが2個，赤カードが3個入っている。はじめにAの袋から同時に2個のカードを取り出し，続いて，Bの袋から同時に2個のカードを取り出す。

\vspace{1.5em}

(1) \quad Aから2個の白カードを取り出し，Bからは白カードと赤カードをそれぞれ1個ずつ取り出す確率は$\dfrac{\ab{\color{red}{\textsf{J}}}}{\ab{\color{red}{\textsf{KL}}}}$ \quad 

\vspace{1.5em}

(2) \quad 取り出した4個のカードの中に，3個の白カードと1個の赤カードが入っている確率は \quad $\dfrac{\ab{\color{red}{\textsf{M}}}}{\ab{\color{red}{\textsf{N}}}}$ \quad である。

\vspace{1.5em}

(3) \quad 取り出した4個のカードがすべて同じ色である確率は \quad $\dfrac{\ab{\color{red}{\textsf{O}}}}{\ab{\color{red}{\textsf{PQ}}}}$ \quad である。

\vspace{1.5em}

(4) \quad 取り出した4個のカードの中に含まれる白カードが2個以下である確率は \quad $\dfrac{\ab{\color{red}{\textsf{RS}}}}{\ab{\color{red}{\textsf{TU}}}}$ \quad である。\\

\newpage
\begin{center}
- 計算欄 (memo) -
\end{center}
\newpage
%第二部分
\noindent
\begin{tikzpicture}
\node[draw, rectangle, minimum width=1cm, minimum height=0.8cm, font=\Huge] {II};
\end{tikzpicture}
\\
\\
正の数からなる数列 $a_1, a_2, a_3, \cdots$ は

\vspace{0.5em}

\begin{align*}
a_1 &= 1, \quad a_2 = 2 \\[0.5em]
(a_n)^2 a_{n-2} &= (a_{n-1})^3 \quad (n = 3, 4, \cdots) \quad \cdots\cdots\cdots (1)
\end{align*}

を満たしている。このとき，$\displaystyle\lim_{n \to \infty} a_n$ を求めよう。

\vspace{1em}

\indent
(1) の両辺の常用対数を考えて

\vspace{0.5em}

\[
\ab{\color{red}{\textsf{A}}} \log_{2} a_n + \log_{2} a_{n-2} = \ab{\color{red}{\textsf{B}}} \log_{2} a_{n-1}
\]

\vspace{0.3em}

を得る。いま，$b_n = \log_{2} a_n \quad (n = 1, 2, \cdots)$ とおくと，この式は

\vspace{0.5em}

\[
\abb{\color{red}{\textsf{A}}} b_n + b_{n-2} = \abb{\color{red}{\textsf{B}}} b_{n-1} \quad \cdots\cdots\cdots (2)
\]

\vspace{0.3em}

となる。(2) を変形すると

\vspace{0.5em}

\[
b_n - b_{n-1} = \frac{1}{\ab{\color{red}{\textsf{C}}}} (b_{n-1} - b_{n-2}) \quad (n = 3, 4, \cdots)
\]

\vspace{0.3em}

となるから

\vspace{0.5em}

\[
b_n - b_{n-1} = \left( \frac{1}{\abb{\color{red}{\textsf{C}}}} \right)^{n-\sab{\color{red}{\textsf{D}}}} (b_2 - b_1) \quad (n = 2, 3, \cdots) \quad \cdots\cdots\cdots (3)
\]

\vspace{0.3em}

が成り立つ。\\[0.5em]
ここで，$b_1 = \ab{\color{red}{\textsf{E}}}$，$b_2 = \ab{\color{red}{\textsf{F}}}$ であるから，(3) より

\vspace{0.5em}

\[
b_n = \sum_{k=2}^{n} \left( \frac{1}{\abb{\color{red}{\textsf{C}}}} \right)^{k-\sab{\color{red}{\textsf{G}}}}
\]

\vspace{0.3em}

を得る。よって

\vspace{0.5em}

\[
b_n = \ab{\color{red}{\textsf{H}}} - \left( \frac{1}{\abb{\color{red}{\textsf{C}}}} \right)^{n-\sab{\color{red}{\textsf{I}}}}
\]

\vspace{0.3em}

である。したがって

\vspace{0.5em}

\[
\lim_{n \to \infty} a_n = \ab{\color{red}{\textsf{J}}}
\]

\vspace{0.3em}

である。

% \textbf{問2}\qquad 2つの2次関数

% \[
%   f(x) = \frac{1}{3}x^2 + ax - b, \qquad g(x) = -x^2 + cx + b
% \]

% が、次の2つの条件 (A), (B) を満たすような $a$, $b$, $c$ を求めよう。

% \begin{itemize}
%   \item[(A)] グラフ $y=f(x)$ とグラフ $y=g(x)$ は2つの直線 $x=-1,~x=3$ 上で交わっている。\\[0.5em]
%   \item[(B)] $g(x)$ の最大値と $f(x)$ の最小値の差は $\dfrac{16}{3}$ である。
% \end{itemize}
% \ \\[0.5em]
% 条件(A) より2つの2次関数のグラフは直線 $x=-1$ 上で交わっているから
% \[
%   3a + \ab{\color{red}{\textsf{M}}} b - \ab{\color{red}{\textsf{N}}} c = \ab{\color{red}{\textsf{O}}} \qquad \cdots\ (1)
% \]
% であり、また、直線 $x = 3$ 上で交わっているから
% \[
%   \ab{\color{red}{\textsf{P}}} a - 2b - \ab{\color{red}{\textsf{Q}}} c = -12 \qquad \cdots\ (2)
% \]
% を得る。①と②より $b = \ab{\color{red}{\textsf{R}}}$ である。\\[0.5em]

% \noindent
% 次に、条件(B) と $b = \abb{\color{red}{\textsf{R}}}$ より\\[0.5em]
% \[
%   \ab{\color{red}{\textsf{S}}} a^2 + 3c^2 = \ab{\color{red}{\textsf{TU}}}
% \]
% を得る。これらから\\[0.5em]
% \[
%   a = -\frac{\ab{\color{red}{\textsf{V}}}}{\ab{\color{red}{\textsf{W}}}}, \qquad
%   c = \ab{\color{red}{\textsf{X}}}
% \]
% となる。

\newpage
\begin{center}
- 計算欄 (memo) -
\end{center}
\newpage
%第三部分
\noindent
\begin{tikzpicture}
\node[draw, rectangle, minimum width=1cm, minimum height=0.8cm, font=\Huge] {III};
\end{tikzpicture}
\\

$x$ が不等式

\vspace{0.5em}

\[
2(\log_{\frac{1}{3}} x)^2 + 9\log_{\frac{1}{3}} x + 9 \leq 0 \quad \cdots\cdots\cdots (1)
\]

\vspace{0.3em}

を満たすとき，関数

\vspace{0.5em}

\[
f(x) = (\log_3 x)\left(\log_3 \frac{x}{3}\right)\left(\log_3 \frac{x}{27}\right) \quad \cdots\cdots\cdots (2)
\]

\vspace{0.3em}

の最大値を求めよう。

\vspace{1em}

(1) を満たす $x$ の値の範囲は

\vspace{0.5em}

\[
\ab{\color{red}{\textsf{A}}}\sqrt{\ab{\color{red}{\textsf{B}}}}\leq x \leq \ab{\color{red}{\textsf{CD}}}
\]

\vspace{0.3em}

である。

ここで，$\log_3 x = t$ とおくと，$t$ のとる値の範囲は

\vspace{0.5em}

\[
\cfrac{\ab{\color{red}{\textsf{E}}}}{\ab{\color{red}{\textsf{F}}}} \leq t \leq \ab{\color{red}{\textsf{G}}}
\]

\vspace{0.3em}

である。

また，(2) の右辺を $t$ で表して，その式が表す関数を $g(t)$ とおくと，その導関数は

\vspace{0.5em}

\[
g'(t) = \ab{\color{red}{\textsf{H}}} t^2 - \ab{\color{red}{\textsf{I}}} t + \ab{\color{red}{\textsf{J}}}
\]

\vspace{0.3em}

である。したがって，$f(x)$ は $x = \ab{\color{red}{\textsf{KL}}}$ で最大値 $\ab{\color{red}{\textsf{M}}}$ をとる。


\newpage
\begin{center}
- 計算欄 (memo) -
\end{center}
\newpage
%第四部分
\noindent
\begin{tikzpicture}
\node[draw, rectangle, minimum width=1cm, minimum height=0.8cm, font=\Huge] {IV};
\end{tikzpicture}
\\

\noindent
$a > 0$ とする。曲線 $y = \sqrt{x}e^{-x}$ と $x$ 軸および $x$ 軸上の点 $\mathrm{A}(a, 0)$ を通る直線 $x = a$ で囲まれた部分を，$x$ 軸の周りに 1 回転してできる立体の体積を $V$ とする。

\vspace{1em}

(1) \quad $V$ は $a$ の関数として

\vspace{0.5em}

\[
V = -\frac{\pi}{4}\left\{\left(\ab{\color{red}{\textsf{A}}}a + \ab{\color{red}{\textsf{B}}}\right)e^{-\sab{\color{red}{\textsf{C}}}a} - \ab{\color{red}{\textsf{D}}}\right\}
\]

\vspace{0.3em}

と表される。

\vspace{1em}

(2) \quad 点 A は原点を出発して，$x$ 軸上を正の方向に移動し，その $t$ 秒後の速度を $4t$ とする。(すなわち、$\dfrac{da}{dt} = 4t$)このとき，$t$ 秒後の $V$ の変化率を求めると

\vspace{0.5em}

\[
\dfrac{dV}{dt} = \ab{\color{red}{\textsf{E}}}\pi t^{\sab{\color{red}{\textsf{F}}}}e^{-\sab{\color{red}{\textsf{G}}}t^{\sab{\color{red}{\textsf{H}}}}}
\]

\vspace{0.3em}

である。また，この変化率が最も大きくなるのは

\vspace{0.5em}

\[
t = \frac{\sqrt{\ab{\color{red}{\textsf{I}}}}}{4}
\]

\vspace{0.3em}

のときで，そのときの $V$ の値は

\vspace{0.5em}

\[
V = -\frac{\pi}{8}\left(\ab{\color{red}{\textsf{J}}}\ e^{-\frac{\sab{\color{red}{\textsf{K}}}}{\sab{\color{red}{\textsf{L}}}}} - \ab{\color{red}{\textsf{M}}}\right)
\]

\vspace{0.3em}

である。


\newpage
\begin{center}
- 計算欄 (memo) -
\end{center}

\end{CJK}

\end{document}